% GENERATED FILE. Edit canonical lessons, then run npm run generate:books.
% canonical-source-hash: fnv1a64:55302e018539aa2c
% canonical-lessons: IT-C104-gonna, IT-C104-cappotto, IT-C104-scarpe, IT-C104-calzini, IT-C104-camicetta

\chapter{Skirt, Coat, Shoes, Socks, Blouse}
\label{ch:it-skirt-coat-shoes-socks-blouse}
\hlchaptermodality{104}

\begin{chapteropening}
\textbf{By the end of this chapter:} \emph{I can say a skirt, a coat, shoes, socks and a blouse.}

Complete the last lesson of chapter 104: I can say a skirt, a coat, shoes, socks and a blouse.
\end{chapteropening}

\section[la gonna]{la gonna — a skirt}

\label{lesson:IT-C104-gonna}



\begin{quote}
Before the new one: say the Italian for a scarf, then the Italian for boots.
\end{quote}

\subsection*{You'll want to know: la gonna}
\textbf{la gonna} — \textquotedblleft{}a skirt\textquotedblright{}.

\begin{cousinweb}
\textbf{gonna} $\leftarrow$ Late Latin \emph{gunna}, a garment of fur or leather — inherited. (EtymDB 2.1) Where \emph{gunna} itself came from is uncertain.
\end{cousinweb}

\begin{grammarlens}[title={What you've built}]
One more word for this chapter.
\end{grammarlens}

\subsection*{Your turn}
\begin{itemize}
\raggedright
  \item \emph{Say it:} \emph{la gonna}
  \item \emph{Say it:} \emph{la gonna}, once more
  \item \emph{Say it:} say \emph{la gonna}
  \item \emph{Recall:} say the Italian for a scarf, then the Italian for boots, then say \emph{la gonna} again
\end{itemize}

\begin{tcolorbox}[breakable,colback=teal!4,colframe=teal!35!black,title={Before you move on}]
What does la gonna mean? (\textquotedblleft{}A skirt\textquotedblright{}.) Say it once more.
\end{tcolorbox}
\section[il cappotto]{il cappotto — a coat}

\label{lesson:IT-C104-cappotto}



\begin{quote}
Before the new one: say the Italian for boots, then the Italian for a skirt.
\end{quote}

\subsection*{You'll want to know: il cappotto}
\textbf{il cappotto} — \textquotedblleft{}a coat\textquotedblright{}.

\begin{cousinweb}
\textbf{cappotto} — built from an older word for a hooded cloak, which $\leftarrow$ Late Latin \emph{cappa}. (EtymDB 2.1)
\end{cousinweb}

\begin{grammarlens}[title={What you've built}]
One more word for this chapter.
\end{grammarlens}

\subsection*{Your turn}
\begin{itemize}
\raggedright
  \item \emph{Say it:} \emph{il cappotto}
  \item \emph{Say it:} \emph{il cappotto}, once more
  \item \emph{Say it:} say \emph{il cappotto}
  \item \emph{Recall:} say the Italian for boots, then the Italian for a skirt, then say \emph{il cappotto} again
\end{itemize}

\begin{tcolorbox}[breakable,colback=teal!4,colframe=teal!35!black,title={Before you move on}]
What does il cappotto mean? (\textquotedblleft{}A coat\textquotedblright{}.) Say it once more.
\end{tcolorbox}
\section[le scarpe]{le scarpe — shoes}

\label{lesson:IT-C104-scarpe}



\begin{quote}
Before the new one: say the Italian for a skirt, then the Italian for a coat.
\end{quote}

\subsection*{You'll want to know: le scarpe}
\textbf{le scarpe} — \textquotedblleft{}shoes\textquotedblright{}.

\begin{cousinweb}
\textbf{scarpe} — its origin is not recorded in EtymDB 2.1.
\end{cousinweb}

\begin{grammarlens}[title={What you've built}]
One more word for this chapter.
\end{grammarlens}

\subsection*{Your turn}
\begin{itemize}
\raggedright
  \item \emph{Say it:} \emph{le scarpe}
  \item \emph{Say it:} \emph{le scarpe}, once more
  \item \emph{Say it:} say \emph{le scarpe}
  \item \emph{Recall:} say the Italian for a skirt, then the Italian for a coat, then say \emph{le scarpe} again
\end{itemize}

\begin{tcolorbox}[breakable,colback=teal!4,colframe=teal!35!black,title={Before you move on}]
What does le scarpe mean? (\textquotedblleft{}Shoes\textquotedblright{}.) Say it once more.
\end{tcolorbox}
\section[i calzini]{i calzini — socks}

\label{lesson:IT-C104-calzini}



\begin{quote}
Before the new one: say the Italian for a coat, then the Italian for shoes.
\end{quote}

\subsection*{You'll want to know: i calzini}
\textbf{i calzini} — \textquotedblleft{}socks\textquotedblright{}.

\begin{cousinweb}
\textbf{calzini} — a small form (\emph{-ino}) of the word for a stocking, which $\leftarrow$ Vulgar Latin \emph{calcea}, a covering for foot and leg. (EtymDB 2.1)
\end{cousinweb}

\begin{grammarlens}[title={What you've built}]
One more word for this chapter.
\end{grammarlens}

\subsection*{Your turn}
\begin{itemize}
\raggedright
  \item \emph{Say it:} \emph{i calzini}
  \item \emph{Say it:} \emph{i calzini}, once more
  \item \emph{Say it:} say \emph{i calzini}
  \item \emph{Recall:} say the Italian for a coat, then the Italian for shoes, then say \emph{i calzini} again
\end{itemize}

\begin{tcolorbox}[breakable,colback=teal!4,colframe=teal!35!black,title={Before you move on}]
What does i calzini mean? (\textquotedblleft{}Socks\textquotedblright{}.) Say it once more.
\end{tcolorbox}
\section[la camicetta]{la camicetta — a blouse}

\label{lesson:IT-C104-camicetta}



\begin{quote}
Before the new one: say the Italian for shoes, then the Italian for socks.
\end{quote}

\subsection*{You'll want to know: la camicetta}
\textbf{la camicetta} — \textquotedblleft{}a blouse\textquotedblright{}.

\begin{cousinweb}
\textbf{camicetta} — built from the word for \textquotedblleft{}shirt\textquotedblright{} with \emph{-etta}, \textquotedblleft{}little\textquotedblright{}. EtymDB 2.1 has no entry for it; the parts are plain to see. The shirt word $\leftarrow$ Late Latin \emph{camisia} (EtymDB 2.1).
\end{cousinweb}

\begin{grammarlens}[title={What you've built}]
One more word for this chapter.
\end{grammarlens}

\subsection*{Your turn}
\begin{itemize}
\raggedright
  \item \emph{Say it:} \emph{la camicetta}
  \item \emph{Say it:} \emph{la camicetta}, once more
  \item \emph{Say it:} say \emph{la camicetta}
  \item \emph{Recall:} say the Italian for shoes, then the Italian for socks, then say \emph{la camicetta} again
\end{itemize}

\begin{tcolorbox}[breakable,colback=teal!4,colframe=teal!35!black,title={Before you move on}]
What does la camicetta mean? (\textquotedblleft{}A blouse\textquotedblright{}.) Say it once more.
\end{tcolorbox}
